% !TeX root = ../main.tex
%------------------------------------------------------------------------------
%  Discontinuities and Structural Analysis
%  Sources: Hoek & Brown (1980) Ch. 4 and Ch. 7; Brady & Brown (2006) Ch. 3
%  and Ch. 9; Goodman & Shi (1985) block theory.
%  Net construction procedures are in Appendix~\ref{app:stereonet}.
%------------------------------------------------------------------------------
\chapter{Discontinuities and Structural Analysis}
\label{ch:structural-analysis}

Structurally controlled failure is one of the three principal underground
failure modes, alongside stress-induced failure of the rock material and
large deformation of weak ground (Chapter~\ref{ch:empirical-design}), and it
is the one least often anticipated. A rock mass rated
good by any classification system can release a wedge from the crown if three
discontinuity sets intersect above the opening in an adverse
geometry, whatever the intact strength. The mechanism is
geometric before it is mechanical: gravity acts on a block that the
discontinuities have already defined, and the design question is whether that
block can move.

Stereographic projection is the tool for answering that question. The technique
is usually introduced through slope stability, but it applies equally
underground: roof wedges, sidewall blocks, portal failures and
intersection collapses are all identified by the same construction. The
underground problem is in some ways the harder one, because a tunnel presents
several differently oriented free surfaces simultaneously and a set that is
harmless in the crown may be critical in a sidewall.

This chapter establishes the technique and applies it underground. The
slope kinematic tests are treated only so far as the underground and portal
cases require (Section~\ref{sec:mode-and-orientation}); the step-by-step net
constructions are collected in Appendix~\ref{app:stereonet}.

\section{Projection and discontinuity data}
\label{sec:projection-and-data}

\subsection{Projection fundamentals}
\label{sec:projection-fundamentals}

Stereographic projection represents three-dimensional orientations on a plane.
A plane through the centre of a reference sphere intersects the lower
hemisphere in a great circle; a line intersects it at a point. Projecting
those from the zenith onto the equatorial plane gives the equal-angle net, on
which a plane appears as a great circle and, equivalently, as a single point,
its pole, which is the projection of the normal. The equal-area net is not a
projection from a point: its radius is rescaled so that equal solid angles on
the sphere occupy equal areas on the plot.

Two projections are therefore in use. The equal-angle, or Wulff, net preserves
angular relationships and is the correct choice for geometric construction
of intersections, true angles and rotations. The equal-area, or Schmidt, net
preserves area and is the correct choice for statistical work. Because the
equal-angle construction stretches area towards the perimeter, a uniform
scatter of poles plots at a higher apparent density near the centre
(Figure~\ref{fig:net-density-bias}), so
contouring pole density on an equal-angle net biases the result towards the
centre. Using the wrong net for the purpose in hand is a common error that the plot does not reveal.

\begin{figure}[htbp]
\centering
\datafigure[\textwidth]{ch09/net-density-bias}
\caption{The same 600 directions, drawn from a uniform distribution over the lower hemisphere, on the two nets ($10^\circ$ grid). The broken circle of half the net radius encloses a quarter of the plot area; it holds 40~per cent of the points on the equal-angle net but 26~per cent, close to the expected 25, on the equal-area net.}
\label{fig:net-density-bias}
\end{figure}

Conventions cause more mistakes than the geometry does. A plane may be
recorded as dip and dip direction or as strike and dip, and the strike
convention itself has right-hand and quadrant variants. Mixing them produces
planes rotated by ninety degrees, which is usually obvious, or dipping the
wrong way, which is not. The convention should be fixed for a project before
any data are collected, stated on every plot, and checked whenever data from
different sources are combined (Section~\ref{sec:net-setup}).

\subsection{Sets, intersections and domains}
\label{sec:sets-and-intersections}

Field data arrive as a scatter of individual measurements, and the first task
is to decide how many sets are present. Pole density contouring identifies
clusters (Figure~\ref{fig:pole-contouring}); the mean orientation of each cluster defines a set, and the
dispersion about that mean (conveniently summarised by the Fisher
concentration~\citep{fisher1953dispersion}, large for a tightly defined set
and approaching zero for a scatter indistinguishable from random)
describes how tightly it is defined. Judgement is required.
Over-interpretation invents sets from noise; under-interpretation merges two
distinct sets into one mean orientation that describes neither, which is worse
because the merged mean may correspond to no plane that exists.

\begin{figure}[htbp]
\centering
\datafigure[0.62\textwidth]{ch09/pole-contouring}
\caption{Pole density on the equal-area net for a synthetic sample of 200 poles: three Fisher-distributed sets about the means marked (J1 $40^\circ/225^\circ$, J2 $65^\circ/225^\circ$, J3 $80^\circ/045^\circ$) and 20 random poles, contoured at 2, 4, 6 and 8~per cent per 1~per cent area (maximum 11.4). J1 and J2 merge at the lowest contour; J3 wraps across the primitive.}
\label{fig:pole-contouring}
\end{figure}

Random discontinuities that belong to no set must still be considered. A single persistent
feature can release a block when combined with two systematic sets,
and the largest wedges are frequently formed this way.

The intersection of two planes is the most important construction for
underground work (Section~\ref{sec:construct-intersection}). It plots as a point, its plunge and trend are read directly,
and it is the direction along which a wedge formed by those two planes will
slide. Most wedge analysis reduces to locating a small number of intersection
lines and comparing their orientation with the excavation geometry.

Structure varies along an alignment, and a single stereonet for a whole
project conceals the variation that governs design. The data should be divided
into structural domains (reaches within which the sets are consistent)
with boundaries drawn at faults, lithological changes or changes in
deformation history. A domain boundary is a geological judgement, and where a
domain contains too few measurements to define its sets, the shortfall is a
gap in the investigation and does not justify merging it with its neighbour.

\subsection{Sampling bias and data quality}
\label{sec:sampling-bias}

Discontinuity sampling is systematically biased. A borehole or scanline
intersects planes perpendicular to it in proportion to their frequency, and
planes parallel to it hardly at all. A vertical borehole therefore
under-samples steeply dipping joints, which in a tunnel are often the
critical ones.

The Terzaghi correction~\citep{terzaghi1965sources} weights each observation by the reciprocal of the
sine of the acute angle $\delta$ between the discontinuity and the sampling
line, $w = 1/\sin\delta$, restoring the expected frequency
(Section~\ref{sec:statistical-treatment-of-orientation-data}). The correction has one limitation: as that angle approaches zero
the weight diverges, so a single measurement of a near-parallel set can be
multiplied into an apparent cluster. A cut-off angle, conventionally
$15^\circ$ to $20^\circ$, is applied to prevent this
(Figure~\ref{fig:terzaghi-weight}), and the
consequence is that sets close to parallel with the borehole remain
under-represented whatever correction is applied. The only remedy
is to sample in more than one direction.

\begin{figure}[htbp]
\centering
\datafigure[0.62\textwidth]{ch09/terzaghi-weight}
\caption{Terzaghi weight $w = 1/\sin\delta$ against the angle $\delta$ between discontinuity and sampling line~\citep{terzaghi1965sources}, capped at cut-offs of $20^\circ$ ($w = 2.9$) and $15^\circ$ ($w = 3.9$). Sets in the shaded band stay under-represented whatever the correction.}
\label{fig:terzaghi-weight}
\end{figure}

Four other errors recur. Magnetic declination is left uncorrected, or is
corrected twice. Data sets from different campaigns are merged without checking
that they used the same convention. Contour plots are drawn from too few poles to
support the contour interval. Most consequentially, a mean set
orientation is carried forward into analysis as though it had no scatter, when
the scatter determines whether an adverse wedge is possible.

\section{Block formation and orientation}
\label{sec:block-formation}

\subsection{Removability}
\label{sec:removability-projection}

Whether a block can move at all is established before its stability is
analysed. A block is of engineering interest only if it is finite (bounded on all sides)
and removable, meaning it can be displaced into the excavation without
interpenetrating the surrounding rock.

Block theory~\citep{goodman1985block} establishes this geometrically. The discontinuity sets bounding a
block define a joint pyramid, the set of directions in which the block is free
to move as determined by its faces; the excavation surface defines an
excavation pyramid, the directions occupied by rock. If the joint pyramid is
not itself empty, so that the faces permit some direction of movement, and
shares no direction with the excavation pyramid, the block is finite and can
be withdrawn into the opening (Figure~\ref{fig:block-removability}). If they
overlap, the block is
constrained by the surrounding rock and cannot move regardless of its weight.

\begin{figure}[htbp]
\centering
\datafigure[\textwidth]{ch09/block-removability}
\caption{Blocks at a roof in section (schematic, after \citet{goodman1985block}). Removable: the faces converge away from the free surface and the block can fall. Tapered: finite, but wider above than at the surface, so the surrounding rock locks it. Infinite: the faces never close and no block is formed.}
\label{fig:block-removability}
\end{figure}

The result depends only on orientation, not on
strength, size or stress. It is decided by the discontinuity sets and the
orientation of the free surface, and it is therefore available as soon as
structural data exist, before any parameters have been assigned. The test
assumes a planar free surface. A curved profile is treated by applying it to
the tangent plane at each location of interest, so the verdict varies around
the profile: a block removable at the crown may be constrained at the
springline, and the choice of tangent planes is part of the analysis.

\subsection{Wedges in a tunnel}
\label{sec:tunnel-wedges}

Three intersecting sets in the crown of a tunnel form a tetrahedral wedge.
Its geometry follows from the three plane orientations and the excavation
dimensions: the apex height is fixed by the intersection geometry, the base by
the excavation surface, and the weight by the volume and unit weight (Section~\ref{sec:block-and-wedge-construction}).
Pattern bolts are effective only where they cross the wedge's bounding planes
and anchor in the stable rock beyond the apex (Figure~\ref{fig:r3d-crown-wedge}).

\begin{figure}[htbp]
\centering
\renderfigure{crown_wedge}{%
  \rlab[west]{0.47,0.66}{0.86,0.64}{Wedge}
  \rlab[west]{0.575,0.78}{0.86,0.80}{Bolts}
  \rlab[east]{0.30,0.47}{0.16,0.62}{Tunnel}
}
\caption{A tetrahedral wedge (carmine) released from the crown of a 6~m tunnel, its apex 2.2~m above the crown, shown dropped by 0.3~m; the rock block is drawn by its edges only. Six 4~m radial bolts on a 0.9 by 3~m pattern (malachite) cross the wedge into the rock beyond (illustrative geometry).}
\label{fig:r3d-crown-wedge}
\end{figure}

Wedge weight scales with the cube of the linear dimension, so doubling the
span multiplies the weight of the geometrically similar wedge by eight. This
is why large-span caverns are structurally demanding in a way that
extrapolation from tunnel experience does not convey, and why span is the
dominant variable in cavern support design
(Chapter~\ref{ch:openings}).

The same three sets form different wedges on each free surface. The crown, the
two sidewalls and the face each present a differently oriented plane, and a set
combination that is harmless in one may be critical in another. An analysis
that considers only the crown has examined one of four free surfaces. Sidewall
wedges are additionally penalised by the low clamping stress on a relaxed
wall, which reduces the frictional resistance available on their sliding
planes.

\subsection{Failure mode and excavation orientation}
\label{sec:mode-and-orientation}

The failure mode follows from the geometry. If the line of intersection
plunges more steeply than the friction angle (Table~\ref{tab:joint}) and daylights into the
excavation, sliding occurs along it: two-plane sliding on the wedge's
bounding planes, the third acting as a release surface. If only one plane
dips more steeply than the friction angle, daylights in the free surface and
has a dip direction close to that of the surface (within about $20^\circ$ is
the usual limit), sliding occurs on that plane alone and the other two
release. If no intersection daylights and the block is in the crown, the mode
is free fall: the faces separate and no shear resistance is mobilised along
them, so that only cohesion across the faces and any clamping stress still
acting can delay detachment. Each mode mobilises a different resistance, and identifying it
correctly is a precondition for the stability calculation. The two sliding
tests, applied with a portal face or cut slope as the free surface, are the
kinematic assessment of a portal (Figure~\ref{fig:portal-kinematics}).

\begin{figure}[htbp]
\centering
\datafigure[\textwidth]{ch09/portal-kinematics}
\caption{Kinematic tests on a portal face, $\phi = 35^\circ$. Left: the intersection I (plunge/trend $41^\circ/150^\circ$) of Ja $60^\circ/090^\circ$ and Jb $60^\circ/210^\circ$ lies between the friction circle (broken) and the $70^\circ/135^\circ$ face, so the wedge can slide. Right, poles on a $55^\circ/220^\circ$ face: J1 $40^\circ/225^\circ$ lies in the planar sliding zone; J2 $65^\circ/225^\circ$ does not daylight.}
\label{fig:portal-kinematics}
\end{figure}

Excavation orientation is the design variable that structural analysis
provides, and it can remove a hazard rather than support it. Rotating the
excavation on the stereonet changes the relationship between the free surfaces
and the discontinuity sets, and can convert a large critical wedge into a
small stable one (Figure~\ref{fig:crown-wedge-orientation};
Section~\ref{sec:construct-rotation}). Where the alignment has freedom, testing several
orientations against the structural data costs almost nothing and may save a
support class over a substantial length.

\begin{figure}[htbp]
\centering
\datafigure[\textwidth]{ch09/crown-wedge-orientation}
\caption{Largest crown wedge of three illustrative sets ($65^\circ/070^\circ$, $60^\circ/280^\circ$, $40^\circ/180^\circ$) in a flat roof of 10~m span, against tunnel axis trend. Right, plan: base triangles at the worst trend ($160^\circ$: 215~m\textsuperscript{3}, apex 6.2~m) and the best ($070^\circ$: 23~m\textsuperscript{3}, apex 2.9~m), broken lines running to the apex. Rotation alone reduces the wedge ninefold.}
\label{fig:crown-wedge-orientation}
\end{figure}

The complication is that the alignment
that minimises structural hazard is rarely the one that minimises boundary
stress concentration (Chapter~\ref{ch:in-situ-stress}), and the two must be
traded against each other explicitly, with the governing mechanism identified
rather than assumed.

\section{From geometry to support}
\label{sec:geometry-to-support}

Once the wedge is identified and its mode established, the analysis becomes
mechanical: weight drives; friction, cohesion and clamping stress on the
sliding planes resist; and support makes up the deficit, its optimum
orientation generally inclined to the sliding direction rather than normal to
the surface. The limit-equilibrium treatment of an identified block,
including the effect of water pressure on the sliding planes, is the subject
of Section~\ref{sec:blocky-rock}; the calculation sequence is set out in
Appendix~\ref{app:calculations}, and the design of the reinforcement itself in
Chapter~\ref{ch:support}.

Orientation data are scattered, and treating a mean set orientation as exact
produces one wedge where a distribution of wedges exists. Sampling orientations
from their observed distributions and repeating the construction yields a
distribution of wedge sizes and factors of safety, which is a
more informative design input because the critical case usually lies
in the tail of the distribution rather than at the mean. Structural
analysis becomes probabilistic at this point, and the methods are those of
Chapter~\ref{ch:integrated-design}.

\section{Software and digital workflows}
\label{sec:structural-software}

Block analysis software performs these constructions reliably and quickly:
identifying the critical wedge among all
combinations of several sets on several free surfaces is tedious by hand and
immediate by machine. The risk is that the software will return a
factor of safety, to three decimal places, for a wedge that could not
physically form, from discontinuity sets that were mis-assigned, or at an
orientation that reflects a convention error. Every
result should be checkable against a hand construction, and the constructions of
Appendix~\ref{app:stereonet} are provided for that purpose. Pole and great-circle plotting and the
intersection of two planes are implemented in \companionname\
(Section~\ref{sec:structural-scripts}), whose output can be checked in the
same way; it analyses wedges on a slope or portal face but not blocks in the
roof and walls of an opening, so the sweep over free surfaces and excavation
orientations underground is left to dedicated software or to the hand
constructions.

Data capture has changed more than analysis. Photogrammetry and terrestrial
LiDAR extract discontinuity orientations from exposures and tunnel faces
automatically, in numbers no manual survey could match and from surfaces that
would be unsafe to approach. The output feeds directly into projection and
domain analysis.

Two limitations apply. Remote methods measure only the surfaces
that are exposed and visible, so they carry the orientation bias of the
exposure: a face perpendicular to a joint set will show it well and a
face parallel to it will not, as with a borehole. They also measure
orientation and trace length well, but roughness, aperture, infilling and wall
strength poorly or not at all, so the parameters that determine shear strength
still require direct access to the rock.

\chapterreferences
